\section*{Unified spine (time, selection, updates, and resolution)}
\addcontentsline{toc}{section}{Unified spine (time, selection, updates, and resolution)}
\label{sec:unified_spine}

\paragraph{Purpose.}
\AuditTag This reader-facing section compresses a recurring set of structural themes into a single vocabulary:
(i) time as tick and as executed update order,
(ii) deterministic CAP closure as finite-family selection with explicit tie-breaks,
(iii) updates and branching as changes of admissible candidates/objectives under fixed constraints,
and (iv) resolution flow (uplift/RG/forced zoom) as an update trajectory of protocol states.
None of this section is used as a premise in theorem-level folding proofs; it is an organizational index.

\subsection*{Three time notions and their layer placement}
\AuditTag We distinguish three time notions throughout:
\begin{itemize}
\item \textbf{tick time} $t_{\mathrm{tick}}$: the scan iteration index (Definition~\ref{def:tick});
\item \textbf{update time} $t_{\mathrm{upd}}$: the executed order of protocol updates (e.g.\ CAP closures or objective updates), indexed discretely and used to state monotonicity certificates (Section~\ref{subsec:wish_motive_generic_dynamics});
\item \textbf{physical time} $t_{\mathrm{phys}}$: any mapping from tick/update indices to SI units (seconds) or laboratory clock readings, which is always a matching-layer calibration (Remark~\ref{rem:tick_velocity_to_c} and Section~\ref{subsec:arrow_time_calibration_constants}).
\end{itemize}
\AuditTag The first two are intrinsic to the protocol language; the third is a dictionary layer.

\subsection*{Finite-family selection as the universal closure form}
\AuditTag The CAP rule is implemented throughout as deterministic selection on explicit finite candidate families.
Abstractly, a selection problem is a triple $(\mathcal{F},J,\prec)$ consisting of a finite candidate family $\mathcal{F}$,
an explicit objective functional $J:\mathcal{F}\to\RR$, and a deterministic tie-break order $\prec$.
The corresponding deterministic selector $\mathsf{Sel}(\mathcal{F},J,\prec)$ is formalized in Appendix~\ref{app:branching_selection_rigidity}
and instantiated in audit form across the tick\,+\,CAP derivation spine (Appendix~\ref{app:tick_cap_derivation}).

\paragraph{Gaps and robustness.}
\AuditTag When a quantitative objective gap is present, bounded perturbations cannot change the selected minimizer
(gap-stability; Appendix~\ref{app:branching_selection_rigidity} and the RB-D template in Appendix~\ref{app:rigidity_bridge_checklist}).
This is the paper's normalized way to state that a closure is robust rather than an accidental alignment.

\subsection*{Updates, branching, and prefix-consistency}
\AuditTag What is informally called ``branching'' is expressed in finite mathematics by prefix projections and extension fibers:
distinct continuations may share an audited prefix while differing in admissible extensions.
This is formalized by extension fibers $\mathrm{Ext}_m(u)$ under prefix projections $\pi_{m\to m_0}$ (Appendix~\ref{app:branching_selection_rigidity}).
\AuditTag Updates are modeled as maps $(\mathcal{F},J,\prec)\mapsto(\mathcal{F}',J',\prec')$; a ``switch'' occurs when the selected minimizer changes.
If a prefix constraint is preserved, any switch is confined to the corresponding extension fiber (Proposition~\ref{prop:prefix_consistent_switching_constraint}).

\subsection*{Where the HPA-integrated modules sit in the Sel/Upd/gap vocabulary (reader map)}
\label{subsec:sel_upd_gap_hpa_map}
\AuditTag Several companion-HPA modules are integrated into the main text in a way that preserves the same audit discipline.
Their placement in the unified Sel/Upd/gap vocabulary is as follows:
\begin{itemize}
\item \textbf{Embedding and projection readout (Sel + gap).}
The projection readout in Section~\ref{sec:embedding_gap_projection} is a concrete instance of deterministic finite-family selection:
for a fixed synthesized $v\in\CC$, the nearest-point rule $\pi(v)$ is a selector on an explicitly finite candidate subset (proved finite in that section), and the objective gap is the separation between the best and second-best distances.
The cohomology-based ``gap field'' viewpoint provides an additional \AuditTag failure mode not reducible to a scalar objective gap (a nontrivial $H^1$ class obstructs global gauging-away of residuals).
\item \textbf{Cut-and-project bridge (prefix/fiber carrier).}
Section~\ref{sec:cut_and_project_bridge} is the geometric carrier behind the prefix/fiber language:
window words induce staircase prefixes and therefore a canonical prefix projection hierarchy; model sets $\Lambda(W)$ provide a point-set realization of the same branching structure.
\item \textbf{Time-advance dispersion channel (dictionary update).}
Section~\ref{sec:dispersion_time_advance} records a kinematic signature once a specific energy/momentum readout is declared (Assumption~\ref{ass:dispersion_energy_momentum_from_scan}).
Interpreting this signature as a laboratory observable is a \MatchTag dictionary update and must be recorded as such, separate from overhead-induced delay channels.
\item \textbf{Octonionic internal lift and adelic prime-orbit modules (candidate-family updates).}
Sections~\ref{sec:octonion_g2_internal_phase} and \ref{sec:adelic_prime_orbit_module} are recorded as \AuditTag/\InterfaceTag candidate modules:
invoking them in any closure requires declaring an explicit finite family (of lifts/orbits/representatives) and a deterministic tie-break, so that the invocation itself is an instance of Sel/Upd under the IC-13 gate.
\end{itemize}

\subsection*{Resolution flow as protocol-state updates}
\AuditTag Many interface outputs depend on a protocol state $(m,n,K)$ (window length, addressing order, readout kernel),
which is closed by CAP within explicit finite families (Section~\ref{sec:paper_contract} and Appendix~\ref{app:tick_cap_derivation}).
In this paper, three recurring narratives are treated as instances of the same update viewpoint:
\begin{itemize}
\item \textbf{uplift}: changing the window length $m$ and selecting fiber representatives under refinement (Appendix~\ref{app:mass_flow_under_uplift});
\item \textbf{RG}: describing running in the resolution coordinate $r(\mu)$ with explicit scheme/threshold boundaries (Appendix~\ref{app:running_couplings_resolution_flow});
\item \textbf{forced zoom}: engineered threshold crossing that changes the effective admissible protocol budget (Section~\ref{subsec:active_renormalization_synthesis}).
\end{itemize}
\AuditTag The mathematical-layer carrier for these statements is always finite (candidate families, prefix fibers, and deterministic keys);
any calibration to laboratory units remains at the matching layer.

